\section{Основни појмови}
У овом поглављу ћемо да објаснимо основне елементе нашег примера из претходне секције.
Концептуално идеја ТЛА+ је да систем представи као неки граф могућих стања који одређеује
која су почетна, крајња стања и како се прелази у следеће стање система на основу тренутног и правила прелаза.
Ово нас доста подсећа на аутомате из ранијих курсева за комапјлере и лексичку анализу.

\begin{figure}[H]
    \centering
    \begin{tikzpicture}[
        node distance=2.2cm,
        every state/.style={circle, draw, minimum size=1cm},
        ->, >=Stealth, thick,
        lbl/.style={font=\scriptsize, fill=white, inner sep=2pt, align=center}
    ]
        \node[state, initial, initial text={}] (s1)  {$1$};
        \node[state, right=of s1]              (s2)  {$2$};
        \node[state, right=of s2]              (s3)  {$3$};
        \node[right=1.2cm of s3, font=\Large]  (dots){$\cdots$};
        \node[state, right=1.2cm of dots, accepting] (s10) {$10$};

        \draw (s1)   edge node[lbl, above] {$c < 10$\\$c' = c+1$} (s2);
        \draw (s2)   edge node[lbl, above] {$c < 10$\\$c' = c+1$} (s3);
        \draw (s3)   edge node[lbl, above] {$c < 10$\\$c' = c+1$} (dots);
        \draw (dots) edge node[lbl, above] {$c < 10$\\$c' = c+1$} (s10);
        \draw (s10)  edge[loop right] node[lbl, right] {$c \geq 10$\\$c' = c$} (s10);
    \end{tikzpicture}
    \caption{Граф прелаза система Counter: стања су вредности бројача ($c$), прелази одговарају повећавању за 1 све до вредности 10}
\end{figure}
Претходна слика представља граф нашег система  из уводног примера са свим својим могућим стањима. Присећамо се да
почетно стање има само улазну стрелицу, а крајње обележавамо са два концентрична круга. Као што можемо видети почетно стање
нашег система је стање \textbf{1}, а крајње стање је \textbf{10}. Прелаз између стања је могућ ако променљива \textbf{counter}
има вредност мању или једнаку од 10 и ако смо у стањима 1 ди 9. На основу спецификације у тла, тла систем генерише граф стања
система и обилази га и покушава да нађе неке грешке: недоступа стања, бесконачне петље итд.

\subsection{Основни делови tla спецификације}
Да се подсетимо претходног уводног примера и конфигурационог фајла.
\begin{lstlisting}
---- MODULE Counter ----
EXTENDS Naturals

VARIABLE counter

Init ==
    counter = 1

Next ==
    IF counter < 10
    THEN counter' = counter + 1
    ELSE counter' = counter

Spec ==
    Init /\ [][Next]_counter
====
\end{lstlisting}
\begin{lstlisting}
CONSTANTS
    greeting = "Hello"

INIT Init
NEXT Next
\end{lstlisting}
Видимо да уводни пример у потпуности одговара графу који смо нацртали. За велике и озбиљне системе
ручно проверавање стања је практично немогуће и ту нам треба помоћ рачунара. Укратко смо представили
поједностављену верзију како tla систем функционише.
Сваки tla спецификација почиње са кључном речи \textbf{Module}. Симболи \textbf{-\;-\;-}  су део стандардне
синтаксе који означавају име и почетак спецификације. Такође \textbf{=\;=\;=} на самом крају фајла
означава крај спецификације и нема неки већи значај-
\begin{lstlisting}
---- MODULE Counter ----
...
====
\end{lstlisting}

Кључна реч \textbf{EXTENDS Naturals} нам увози библиотеку која нам омогућава да радимо са стандардним аритметичким
операцијама $+, \cdot, >, <$ итд.
\begin{lstlisting}
---- MODULE Counter ----
EXTENDS Naturals
...
====
\end{lstlisting}
Да нисмо увезли ову библиотеку не бисмо могли да радимо са изразима попут $counter < 10$ и $counter + 1$.

\textbf{VARIABLE counter} нам декларише променљиву са именом $counter$.
\begin{lstlisting}
---- MODULE Counter ----
...

VARIABLE counter

Init ==
    counter = 1
...
====
\end{lstlisting}
Променљиве су део стања програма и њихова вредност се мења између стања, тј. при прелазу  између стања
добијају вредности. У нашем случају имамо једну променљиву \texttt{counter} и њена вредност се очигледно
мења при промени стања.
\begin{lstlisting}
Init ==
    counter = 1
\end{lstlisting}
Почетна вреднот променљиве у почетном стању је једнака 1.

Следећи део кода нам говори како се прелази у следеће стање програма и томе служи кључна реч \textbf{Next}.
\begin{lstlisting}
---- MODULE Counter ----
...
Next ==
    IF counter < 10
    THEN counter' = counter + 1
    ELSE counter' = counter
...
====
\end{lstlisting}
Структура \textbf{IF ... THEN ... ELSE} има стандардно значење као и у свим програмским језицима, ако је услов испуњен
улазимо у прву грану, а у супротном у \textbf{ELSE} грану. Оно што се разликује од досадашњег искуства је апостроф
над промљивом \texttt{counter'}. Апостроф овде нам служи да кажемо шта је ново стање променљиве \texttt{counter} и како га добијамо. У овом случају
ново стање променљиве је једнаком старом увећано за 1 ако је вредност из претходног стања мања од 10, а у супротоном
остаје непромењена. Када име променљиве пишемо без апострофа то значи да узимамо њену вредност из претходног стања.

Последњи део спецификације нас асоцира на неки логичку конјукцију два услова. Оператор /\textbackslash\; има 
стандардно значење логичке конјукције $\wedge $, па овај део спецификације можемо тумачити да би спецификација
била исправна потребно је да оба услова буду испуњена.
\begin{lstlisting}
---- MODULE Counter ----
...
Spec ==
    Init /\ [][Next]_counter
====
\end{lstlisting}
Систем је свакако почео у стању \textbf{Init} и вредност променљиве јесте 1 на почетку. Други део конјукције користи
важан оператор \textbf{[]}. Оператор \textbf{[]} значи да неки услов треба да буде испуњен \textbf{увек}, то значи
да има временско трајање, па одатле и добијамо назив TLA - temporal logic of actions. Оператор \textbf{[Next]} нам кажемо
да свако следеће стање при прелазу треба да се добија из стања \textbf{Next}, a \texttt{\_counter} значи да стање може да буде 
непромењено при прелазу. Говорним језиком преведено ова реченица значи да стање почиње из стања \textbf{Init} и прелази могу да се дешавају
само према правилу \textbf{Next} или да стање променљиве \texttt{\_count} остане непромењено. Ако би се десило да у систему
може да се деси неки прелаз који није по правилу tla систем ће да провером да то утврди и пријави грешку.

Потпун преглед оператора се налази у поглављу 4.

\subsection{Конфигурациони фајл}
Конфигурациони фајл има једноставну структуру. У нашем малом примеру дефинише неке константе, дефинише шта је следеће стање
и које је име правила прелаза на следеће могуће знање. 
\begin{lstlisting}
CONSTANTS
    greeting = "Hello"

INIT Init
NEXT Next
\end{lstlisting}
На крају поглавља четири има кратак преглед могућих опција за конфигурациони фајл. 

Упозанавње са осталим могућим операторима и идејама tla спецификације ћемо урадити кроз примере. Треба напоменути
да због специфичности временске компоненте оператора њихова правилна употреба није лака и да је најбоље
упознати се са њима кроз мале практичне примере.

\newpage